\documentclass[11pt]{article}
\usepackage[margin=1in]{geometry}
\usepackage{amsmath,amssymb,amsthm,booktabs,hyperref}
\setlength{\emergencystretch}{1em}
\newcommand{\Q}{\mathbb Q}
\newcommand{\R}{\mathbb R}
\DeclareMathOperator{\End}{End}
\DeclareMathOperator{\Aut}{Aut}
\DeclareMathOperator{\Gal}{Gal}
\DeclareMathOperator{\SO}{SO}
\DeclareMathOperator{\SU}{SU}
\DeclareMathOperator{\ST}{ST}
\newtheorem{proposition}{Proposition}
\title{Arithmetic candidates and Euler-polynomial evidence for $(2,0,2)$ motives}
\author{}
\date{Working note, 23 October 2026}
\begin{document}
\maketitle

\paragraph{Scope and status.}
This note classifies the compact candidates in the two non-CM branches,
subject to the stated monodromy and rationality hypotheses. It records rational
quaternion embedding conditions, but does not classify their Galois realizations,
all rational conjugacy classes of embeddings, or the associated surfaces.
The CM cases are included as boundary strata, without an exhaustive arithmetic
component classification. A new Magma pilot verifies the complete Euler factors
of the modular surface $S_1(9)$ at all 25 primes $p\le101$, $p\ne3$.
No other examples have been assigned to component groups in this audit.

\paragraph{Three endomorphism algebras.}
Let $T$ be a rational polarized rank-four weight-two motive over $\Q$ with
Hodge numbers $(2,0,2)$. Work with semisimple realizations and put
$U_\ell=T_\ell(1)$ and $G_\ell=\overline{\rho_{U,\ell}(G_\Q)}^{\rm Zar}$.
The algebra defined without a cycle conjecture is
\[
 A_\ell=\End_{G_\ell^\circ}(U_\ell)
       =\bigcup_{L/\Q\text{ finite}}\End_{G_L}(U_\ell).
\]
Conjugation defines a finite action of $G_\Q$ on $A_\ell$ factoring through
$\pi_0(G_\ell)$. Separately define $A_H=\End_{\rm HS}(T_B)$ and the algebra
$A_{\rm corr}$ of actions on $T$ of algebraic correspondences over $\overline\Q$,
modulo their kernel on cohomology. The latter has a genuine rational Galois action.
Its finite-dimensionality and the fields of definition of a finite basis imply
that this action factors through a finite extension.

The identifications $A_{\rm corr}=A_H$ and
$A_{\rm corr}\otimes\Q_\ell=A_\ell$ are not automatic. In the surface setting,
the relevant Hodge and Tate questions concern the $T^\vee\otimes T$ summand of
$H^4(X\times X)(2)$, hence codimension-two correspondences. Maximal Picard rank
and the Lefschetz $(1,1)$ theorem on $X$ do not settle them. By contrast, for
$H^1$ of an abelian variety, Hodge endomorphisms are algebraic and Faltings supplies
the arithmetic comparison \cite{faltings}.

The Mumford--Tate conjecture predicts
$G_\ell^\circ=\operatorname{Hg}(T_B)\otimes\Q_\ell$, and hence equality of the
corresponding centralizers under comparison. Rational Galois descent of the
Hodge endomorphisms should still be stated explicitly rather than presumed.

\paragraph{Two fields that must not be confused.}
Let $K_{\rm conn,\ell}$ kill the component action and let $K_{\rm end,\ell}$
kill the action on $A_\ell$. Then
$K_{\rm end,\ell}\subseteq K_{\rm conn,\ell}$, but equality need not hold.
If $G_\ell^\circ=\SO_4$ and $G_\ell=\operatorname O_4$, the endomorphism
algebra is scalar, with trivial Galois action, while the component group is $C_2$.
Thus the endomorphism module cannot replace the orientation character
\[
 \eta=\det U_\ell=\det(T_\ell)/\mu_\ell^2.
\]
In the QM branch the component action on the quaternionic centralizer is
faithful, so these two fields agree for a fixed realization. Their common
rational interpretation across $\ell$ belongs to the comparison framework.

\paragraph{Orientation and lattice determinant.}
For a rational geometric compatible system as above, the two-plane
$\operatorname{Fil}^1 U_{\rm dR}$ is maximal isotropic. Thus the de Rham form is
hyperbolic and has square determinant. The local orthogonal comparison formula
$w_1(q_{\rm dR})=w_1(q_B)+\det\rho_{U,\ell}|_{G_{\Q_\ell}}$
\cite[Theorem 1.3]{cnm}, together with independence of $\ell$, gives
$\eta=\chi_{\det(q_B)}$. In the nonsquare case set
$F=\Q(\sqrt{\det(q_B)})$. It is real. All discriminants here are determinant
squareclasses; the signed discriminant agrees in rank four.

\paragraph{Non-CM compact candidates.}
Totaro's Hodge classification \cite{totaro}, with the connected-monodromy
comparison, leads to the generic $\SO(4)$ and quaternionic $\SU(2)$ connected
groups. In the generic branch the nonsquare determinant forces the full compact
group to be $\operatorname O(4)$; its component group is $C_2$ detected by $\eta$.
The proposed modular description is a twisted Asai system of weight $(1,3)$
over $F$. The present draft contains nonsquare-discriminant candidates but this
audit supplies no independent Hilbert modular Euler-factor match.

\begin{proposition}
Let $H=\SU(2)_L\subset\operatorname O(4)$ act on $\mathbb H$ by left
multiplication. Then $N(H)=\SO(4)$ and $N(H)/H\simeq\SO(3)$.
Impose the componentwise FKRS rationality condition and require an involution
with eigenvalues $1,1,-1,-1$. The possible component groups of a compact group
with identity component $H$ are
\[
 C_2,C_4,C_6,D_2,D_3,D_4,D_6,A_4,S_4,\qquad |D_n|=2n.
\]
Each group is the inverse image of its rotation subgroup under
$\SO(4)\to\SO(3)$.
\end{proposition}
\begin{proof}
Use $\SO(4)=(\SU(2)_L\times\SU(2)_R)/\langle(-1,-1)\rangle$.
An orientation-reversing isometry exchanges the two simple factors, so it
does not normalize $H$. On a component represented by $b\in\SU(2)_R$,
the square of the four-dimensional trace has Haar mean $(\operatorname{tr}b)^2$.
For a projective rotation of order $n$ this restricts $n$ to $1,2,3,4,6$,
as in the analogous argument of FKRS \cite[Section 3.4]{fkrs}.
The finite rotation-group classification then gives
$C_1,C_2,C_3,C_4,C_6,D_2,D_3,D_4,D_6,A_4,S_4$.
The only involution in $H$ is $-I$, which has the wrong eigenvalues.
The required involution therefore projects nontrivially, excluding $C_1,C_3$.
Conversely, any projective involution has a lift squaring to $-1$ in
$\SU(2)_R$; pairing it with an element squaring to $-1$ in $\SU(2)_L$ gives
the required involution in $\SO(4)$.
For rationality of the surviving groups, the $H$-invariants in an algebraic
$\mathrm{GL}_4$ representation form an $\SU(2)_R$ representation on which
the center acts trivially. Its character is an integral combination of
even-highest-weight characters, hence a polynomial with integer coefficients
in $(\operatorname{tr}b)^2$. The allowed values make every component average
integral.
\end{proof}

\paragraph{Rational refinement and initial realization-status table.}
Assume a definite rational quaternion algebra $B$ with rational Galois action
has been identified, so $\Gamma\hookrightarrow\Aut_\Q(B)=B^\times/\Q^\times$.
Here $B_2=(-1,-1)_\Q$, and $K$ denotes an imaginary quadratic subfield of $B$.
``Unassigned'' means not assigned by this audit, not a nonexistence claim.
The last column concerns surface examples and their modular identifications;
all nine abstract quaternionic rotation-group constructions are verified by
\texttt{scripts/qm\_candidates.m}.
\begin{center}\small
\begin{tabular}{llll}
\toprule
$\Gamma$ & Condition for an embedding in $B^\times/\Q^\times$ & $B^\Gamma$ & Example status\\
\midrule
$C_2$ & Any definite $B$; choose $K\subset B$ & $K$ & Unassigned\\
$C_4$ & $\Q(i)\hookrightarrow B$ & $\Q(i)$ & Unassigned\\
$C_6$ & $\Q(\sqrt{-3})\hookrightarrow B$ & $\Q(\sqrt{-3})$ & $S_1(9)$ pilot\\
$D_2$ & Any definite $B$ & $\Q$ & Unassigned\\
$D_3$ & $\Q(\sqrt{-3})\hookrightarrow B$ & $\Q$ & Unassigned\\
$D_4$ & $\Q(i)\hookrightarrow B$ & $\Q$ & Unassigned\\
$D_6$ & $\Q(\sqrt{-3})\hookrightarrow B$ & $\Q$ & Unassigned\\
$A_4$ & $B\simeq B_2$ & $\Q$ & Unassigned\\
$S_4$ & $B\simeq B_2$ & $\Q$ & Unassigned\\
\bottomrule
\end{tabular}
\end{center}
These conditions concern existence of a rational projective embedding, not
existence of a motive with that action. They do not enumerate embeddings or
distinguished complex-conjugation classes. An imaginary quadratic field embeds
in $B$ exactly when no place ramified in $B$ splits in that field.

For completeness, the elementary quaternion argument for the table is as follows.
An element of projective order $4$ generates $\Q(i)$, and an element of order
$3$ or $6$ generates $\Q(\sqrt{-3})$. Any pure quaternion gives a projective
involution. Given a quadratic subfield $K$, an element of its normalizer outside
$K^\times$ gives the inversion action and constructs the indicated dihedral
groups. Two perpendicular pure quaternions construct $D_2$ in any definite $B$.
In $A_4$, an order-three element cyclically permutes the axes of its normal
Klein four group. The squares of representatives of these three axes therefore
have the same rational squareclass; if they are $a,b,-ab$, this forces both
$a$ and $b$ to have squareclass $-1$. Hence $B\simeq(-1,-1)_\Q$.
The Hurwitz quaternion constructions give $A_4$, and adjoining $1+i$ gives
$S_4$. The script checks all nine groups via exact rational conjugation matrices.

\paragraph{Exact Frobenius tests.}
For $\eta(p)=1$ write
$P_p(T)=1-A_pT+B_pT^2-p^2A_pT^3+p^4T^4$.
On a QM component the normalized polynomial is
$1-xyT+(x^2+y^2-2)T^2-xyT^3+T^4$.
Writing $r=y^2$, we obtain
\[
 A_p^2=rB_p+r(2-r)p^2,\qquad
 (n,r)=(1,4),(2,0),(3,1),(4,2),(6,3).
\]
For $n=2$ the relation is just $A_p=0$. Several relations can hold at a special
prime, so the checker returns a set of possible orders. The expected densities
of zero trace for the nine rows are respectively
$1/2,1/4,1/6,3/4,1/2,5/8,7/12,1/4,3/8$; the group counts are exact,
whereas a density claim for primes uses the appropriate distribution argument.
The coefficient identities themselves do not require equidistribution.

In either branch the general reciprocal polynomial is
\[
 P_p(T)=1-A_pT+B_pT^2-\eta(p)p^2A_pT^3+\eta(p)p^4T^4.
\]
If $\eta(p)=-1$, then $B_p=0$ and
$P_p(T)=(1-p^2T^2)(1-A_pT+p^2T^2)$.
In particular, the roots $\pm p$ must not be stripped from the characteristic-zero
transcendental factor merely because they become algebraic in the reduction.

\paragraph{CM boundary.}
Keep separate the primitive quartic CM field case, the algebra $K_1\times K_2$
for nonisomorphic rank-two CM pieces, and $M_2(K)$ for two isomorphic pieces.
The first two have compact connected group $U(1)^2$, with normalizer quotient
the dihedral group of order eight. The last has connected scalar $U(1)$ on
$\mathbb C^2$ and normalizer quotient $\SO(3)\times C_2$.
These are boundary normalizer data, not a completed rational arithmetic list.
Record Hecke characters after a suitable extension together with the descent
action and any finite multiplicity representation.

\paragraph{Magma pilot: $S_1(9)$.}
Put $c=t^2(t-1)$, $b=c(t^2-t+1)$, and
\[
 E:\quad y^2+(1-c)xy-by=x^3-bx^2.
\]
This parametrization satisfies the order-nine equation in the supplied draft.
Its singular fibers are $3I_9,2I_3,3I_1$. Their trivial lattice has rank $30$,
equal to $h^{1,1}$, and $b_2=34$; hence its characteristic-zero transcendental
realization has rank four. Let $\chi$ be a primitive character of order six
modulo $9$. The cuspidal modular-symbols space of weight $3$, character $\chi$,
and sign $+$ has dimension one over $\Q(\zeta_6)$, and yields $f$ in the orbit
labelled $9.3.d.a$ in the draft. The script constructs this form without
downloading eigenvalues. For each tested prime it verifies
\[
 P_p(T)=\operatorname{Nm}_{\Q(\zeta_6)/\Q}
       (1-a_p(f)T+\chi(p)p^2T^2),
\]
as well as its tensor expression with $\tau=1\oplus\chi^{-1}$.
The proposed projective Artin image is $C_6$; its splitting field is
$\Q(\zeta_9)$. Every tested polynomial satisfies the corresponding QM relation.
The log records 25 matches for $p\le101$, $p\ne3$, using Magma 2.29-1.
This is a numerical verification, not a new Faltings--Serre proof. The modular
construction provides the route to a separate theoretical identification.
This audit does not determine the rational quaternion algebra of the surface
or exhibit its algebraic correspondences.

\paragraph{Required record for every further example.}
Each record should specify the surface and characteristic-zero projector;
proof/status of $p_g=2$ and maximal Picard rank; the specialization of its
characteristic-zero Neron--Severi representation; $\eta$; the proposed group,
quaternion algebra and endomorphism action (with separate evidence statuses);
the modular/Hecke/Artin data including all twists and normalizations;
an explicit prime bound, excluded primes and reasons; and the path to a
standalone script and output. Keep geometry, monodromy, and modular matching
in separate status columns. Full Euler agreement does not prove algebraicity
of endomorphisms. For elliptic surfaces, Magma's \texttt{LFunction} removes
the trivial lattice, not arbitrary characteristic-zero Mordell--Weil classes
\cite{magma}; the latter require an additional known factor when present.

\begin{thebibliography}{9}
\bibitem{faltings} G. Faltings, \emph{Finiteness theorems for abelian varieties
over number fields}, Arithmetic Geometry (1986), Chapter II.
\bibitem{totaro} B. Totaro, \emph{Hodge structures of type $(n,0,\ldots,0,n)$},
\url{https://www.math.ucla.edu/~totaro/papers/public_html/structure.pdf}.
\bibitem{fkrs} F. Fite, K. Kedlaya, V. Rotger, A. Sutherland,
\emph{Sato--Tate distributions and Galois endomorphism modules in genus 2},
\url{https://arxiv.org/abs/1110.6638}.
\bibitem{cnm} Ph. Cassou-Nogues and B. Morin,
\emph{Tannakian twists of quadratic forms and orthogonal Nori motives},
\url{https://arxiv.org/abs/1511.03022}.
\bibitem{magma} Magma Handbook, \emph{Elliptic curves over function fields},
overview and L-function sections,
\url{https://magma.maths.usyd.edu.au/magma/handbook/text/1574},
\url{https://magma.maths.usyd.edu.au/magma/handbook/text/1581}.
\end{thebibliography}
\end{document}
